% ─────────────────────────────────────────────────────────────────────────────
%  Capítulo 9 – Verificación Formal: Lógica de Hoare y Cálculo WP
% ─────────────────────────────────────────────────────────────────────────────
\chapter{Verificación formal: Lógica de Hoare y cálculo WP}
\label{chap:verificacion}

\section{Fundamentos teóricos}

\subsection{Lógica de Hoare}

La \textbf{Lógica de Hoare}~\cite{hoare1969} es un sistema formal para
razonar sobre la corrección de programas imperativos. Introduce las
\emph{tripletas de Hoare}:
\[
  \hoare{P}{S}{Q}
\]
que se leen: ``si la precondición $P$ vale antes de ejecutar $S$,
entonces la postcondición $Q$ vale tras la ejecución de $S$
(asumiendo que $S$ termina)''.

\begin{definicion}[Corrección parcial]
  $\hoare{P}{S}{Q}$ es \emph{parcialmente correcta} si para toda
  ejecución de $S$ que parte de un estado en el que $P$ vale y
  que termina, el estado final satisface $Q$.
\end{definicion}

\begin{definicion}[Corrección total]
  $\hoare{P}{S}{Q}$ es \emph{totalmente correcta} si además
  toda ejecución de $S$ que parte de un estado en el que $P$ vale
  termina necesariamente.
\end{definicion}

\subsection{Cálculo de Precondiciones Más Débiles}

\textbf{Dijkstra}~\cite{dijkstra1976} introduce el \emph{transformador
de predicados} $\WP{S}{Q}$: la \emph{precondición más débil} de $S$
con respecto a $Q$, es decir, el predicado más débil (el que restringe
menos el estado inicial) tal que:
\[
  \hoare{\WP{S}{Q}}{S}{Q}
\]

La condición de verificación principal para un método anotado es:
\[
  \mathit{VC} \;:\; P \Rightarrow \WP{\mathit{body}}{Q}
\]

Si esta fórmula es \textbf{válida} (lo comprueba Z3 como UNSAT de su negación),
el método es parcialmente correcto.

\section{Reglas del cálculo WP}

\subsection{Regla de la asignación (VER-4)}

\begin{regla}[Asignación]
\[
  \WP{x := e}{Q} \;=\; \subst{Q}{x}{e}
\]
donde $\subst{Q}{x}{e}$ denota $Q$ con cada aparición libre de $x$
sustituida por $e$.
\end{regla}

\begin{ejemplo}
  $\WP{x := x+1}{x > 5} = (x+1) > 5$
\end{ejemplo}

\subsection{Regla de la composición (VER-5)}

\begin{regla}[Composición]
\[
  \WP{S_1;\,S_2}{Q} \;=\; \WP{S_1}{\WP{S_2}{Q}}
\]
\end{regla}

Esto implica procesar la secuencia de sentencias en \emph{orden inverso},
comenzando por la postcondición $Q$ y propagando la precondición hacia
atrás.

\paragraph{Retornos anticipados.}
La regla de composición naive es incorrecta en presencia de \texttt{return}
dentro de ramas \texttt{if}. Para la secuencia:
\[
  \texttt{if}(B)\;\texttt{return}\;e_1;\;\;\texttt{return}\;e_2;
\]
la regla correcta es:
\[
  \WP{\cdot}{Q} \;=\; (B \Rightarrow \subst{Q}{\mathtt{result}}{e_1})
                \;\wedge\; (\neg B \Rightarrow \subst{Q}{\mathtt{result}}{e_2})
\]

La implementación usa \textbf{continuaciones} (CPS): cuando se encuentra
un \texttt{return}, se usa $Q$ directamente (no la continuación de la
composición):

\begin{lstlisting}[style=cpp, caption={wpSeq con continuaciones para retornos anticipados}]
static Predicate* wpSeq(const vector<const StmtNode*>& stmts,
                        int idx, const Predicate* Q) {
    if (idx >= stmts.size()) return Q->clone();
    const StmtNode* s = stmts[idx];
    // return: termina el camino, usa Q directamente
    if (auto* r = dynamic_cast<const ReturnNode*>(s))
        return wpReturn(r, Q);
    // if: bifurca la continuacion
    if (auto* i = dynamic_cast<const IfNode*>(s)) {
        Predicate* B = fromExpr(i->condition.get());
        vector<const StmtNode*> cont(stmts.begin()+idx+1, stmts.end());
        Predicate* wp1 = wpBranchAndCont(i->thenPart.get(), cont, Q);
        Predicate* wp2 = wpBranchAndCont(i->elsePart.get(), cont, Q);
        return new AndPred(new ImpliesPred(B->clone(), wp1),
                           new ImpliesPred(new NotPred(B),  wp2));
    }
    // asignacion: WP(x:=e; resto, Q) = WP(resto, Q)[e/x]
    if (auto* a = dynamic_cast<const AssignNode*>(s)) {
        Predicate* restWP = wpSeq(stmts, idx+1, Q);
        return wpAssign(a, restWP);
    }
    return wpSeq(stmts, idx+1, Q);  // otras: sin efecto logico
}
\end{lstlisting}

\subsection{Regla del condicional}

\begin{regla}[Condicional]
\[
  \WP{\texttt{if}(B)\;S_1\;\texttt{else}\;S_2}{Q}
  \;=\;
  (B \Rightarrow \WP{S_1}{Q}) \;\wedge\; (\neg B \Rightarrow \WP{S_2}{Q})
\]
\end{regla}

\subsection{Regla del retorno}

\begin{regla}[Return]
\[
  \WP{\texttt{return}\;e}{Q} \;=\; \subst{Q}{\mathtt{result}}{e}
\]
\end{regla}

\subsection{Regla del bucle con invariante (VER-6)}

Para la verificación de bucles se requiere una \emph{anotación de invariante}
$I$ suministrada por el programador:

\begin{regla}[Bucle con invariante]
  Para $\texttt{while}(B)\;\{S\}$ con invariante $I$ y postcondición $Q$,
  se generan tres \emph{condiciones de verificación}:
  \begin{align*}
    \mathit{VC}_1 \;(\text{iniciación})\; &: P \Rightarrow \WP{S_{\text{antes}}}{I} \\
    \mathit{VC}_2 \;(\text{preservación})\; &: (I \wedge B) \Rightarrow \WP{S_{\text{body}}}{I} \\
    \mathit{VC}_3 \;(\text{uso})\; &: (I \wedge \neg B) \Rightarrow \WP{S_{\text{después}}}{Q}
  \end{align*}
  y el WP del bucle en sí es $I$ (asumiendo la corrección de $I$).
\end{regla}

La precondición más fuerte al entrar en el bucle
($\WP{S_{\text{antes}}}{I}$) se computa correctamente aplicando
\texttt{wpSeq} a las sentencias que preceden al bucle dentro del bloque.

\section{Jerarquía de predicados}

Los predicados lógicos se representan mediante la clase abstracta
\texttt{Predicate} y sus subclases:

\begin{table}[H]
  \centering
  \caption{Clases de predicados implementadas}
  \label{tab:preds}
  \begin{tabular}{llll}
    \toprule
    \textbf{Clase}    & \textbf{toString()}      & \textbf{toSMT()} & \textbf{Uso} \\
    \midrule
    \texttt{TruePred}     & \texttt{true}        & \texttt{true}   & Postcondición trivial \\
    \texttt{FalsePred}    & \texttt{false}       & \texttt{false}  & VC insatisfacible \\
    \texttt{VarPred(x)}   & \texttt{x}           & \texttt{x}      & Variable de programa \\
    \texttt{ResultPred}   & \texttt{result}      & \texttt{result} & Valor de retorno en @post \\
    \texttt{IntLitPred(n)}& $n$                  & $n$             & Constante entera \\
    \texttt{BinaryPred(op,l,r)} & \texttt{(l op r)} & \texttt{(op l r)} & Aritmética/comparación \\
    \texttt{AndPred(l,r)} & \texttt{(l \&\& r)}  & \texttt{(and l r)} & Conjunción \\
    \texttt{OrPred(l,r)}  & \texttt{(l || r)}    & \texttt{(or l r)}  & Disyunción \\
    \texttt{NotPred(p)}   & \texttt{(!p)}        & \texttt{(not p)}   & Negación \\
    \texttt{ImpliesPred(a,c)} & \texttt{(a ==> c)} & \texttt{(=> a c)} & Implicación \\
    \texttt{NegPred(p)}   & \texttt{(-p)}        & \texttt{(- p)}  & Negación aritmética \\
    \bottomrule
  \end{tabular}
\end{table}

La operación clave es la \textbf{sustitución textual}:
$\texttt{subst}(x, e)$ devuelve una copia del predicado donde cada
\texttt{VarPred(x)} se reemplaza por $e$:

\begin{lstlisting}[style=cpp, caption={Sustitución en BinaryPred}]
Predicate* BinaryPred::subst(const string& var, const Predicate* e) const {
    return new BinaryPred(op, left->subst(var, e), right->subst(var, e));
}
Predicate* VarPred::subst(const string& var, const Predicate* e) const {
    if (name == var) return e->clone();  // sustituir
    return new VarPred(name);           // no afectada
}
\end{lstlisting}

\section{Generación de condiciones de verificación (VER-7)}

El \texttt{VCGenerator} genera cuatro tipos de VCs:

\begin{enumerate}
  \item \textbf{Corrección del método}: $P \Rightarrow \WP{\mathit{body}}{Q}$
  \item \textbf{Invariante de bucle} (iniciación, preservación, uso)
  \item \textbf{Seguridad de memoria} (VER-9): para cada división,
    $P \Rightarrow \mathit{divisor} \neq 0$
  \item \textbf{Terminación} (VER-11): para cada bucle con \texttt{@variant} $V$,
    $(I \wedge B) \Rightarrow V \geq 0$
\end{enumerate}

\section{Exportación a SMT-Lib2 (VER-8)}

Cada VC se exporta como un bloque SMT-Lib2. La negación de la VC se envía
a Z3: si Z3 responde \textbf{UNSAT}, la negación es insatisfacible, por lo
tanto la VC es \textbf{válida} (el programa es correcto en ese punto).

\begin{lstlisting}[style=smt, caption={Formato SMT-Lib2 para una VC}]
(push 1)
; [absValue] correctitud
(declare-fun x () Int)     ; variables locales del metodo
(assert (not              ; negacion de la VC
    (=> true              ; precondicion P
        (and
            (=> (< (- 5) 0) (>= (- (- 5)) 0))
            (=> (not (< (- 5) 0)) (>= (- 5) 0))
        )
    )
))
(check-sat)               ; Z3 responde UNSAT = VC valida
(pop 1)
\end{lstlisting}

\section{Interpretación de resultados (VER-10)}

\begin{table}[H]
  \centering
  \caption{Interpretación de la respuesta de Z3}
  \label{tab:z3resp}
  \begin{tabular}{lll}
    \toprule
    \textbf{Respuesta Z3} & \textbf{Significado}         & \textbf{Reporte} \\
    \midrule
    \texttt{unsat}    & VC válida — programa correcto      & \textcolor{green!50!black}{[OK]} \\
    \texttt{sat}      & VC inválida — contraejemplo existe & \textcolor{red!80!black}{[BUG]} \\
    \texttt{unknown}  & Timeout o lógica demasiado compleja & \textcolor{orange}{[?]} \\
    \bottomrule
  \end{tabular}
\end{table}

\section{Ejemplo completo de verificación}

\subsection{Función \texttt{absValue}}

\begin{lstlisting}[style=java, caption={Método verificado: valor absoluto}]
@pre  true
@post result >= 0
public static int absValue() {
    int x;
    x = -5;
    if (x < 0) return -x;
    return x;
}
\end{lstlisting}

Cálculo WP (con continuaciones):

\begin{align*}
  &\WP{\texttt{return x}}{Q_0} = Q_0[\texttt{x}/\texttt{result}] = x \geq 0 \\
  &\WP{\texttt{if}(x<0)\;\texttt{return}\;{-x};\;\texttt{return x}}{Q_0} \\
  &\quad = (x < 0 \Rightarrow Q_0[{-x}/\texttt{result}])
         \wedge (x \geq 0 \Rightarrow Q_0[\texttt{x}/\texttt{result}]) \\
  &\quad = (x < 0 \Rightarrow -x \geq 0) \wedge (x \geq 0 \Rightarrow x \geq 0) \\
  &\WP{x := -5}{\ldots} = [\text{sustituyendo } x = -5] \\
  &\quad = ((-5) < 0 \Rightarrow 5 \geq 0) \wedge ((-5) \geq 0 \Rightarrow (-5) \geq 0) \\
  &\quad = (\top \Rightarrow \top) \wedge (\bot \Rightarrow \bot) = \top
\end{align*}

VC principal: $\true \Rightarrow \true$ — Z3 responde \texttt{unsat} $\Rightarrow$ \textbf{[OK]}.

\subsection{Función \texttt{wrongMethod} (bug intencional)}

\begin{lstlisting}[style=java, caption={Método con bug intencional}]
@post result > 10
public static int wrongMethod() {
    int x;
    x = 5;
    return x;
}
\end{lstlisting}

$\WP{\texttt{return x}}{Q_1} = Q_1[\texttt{x}/\texttt{result}] = x > 10$

$\WP{x := 5}{x > 10} = 5 > 10 = \false$

VC: $\true \Rightarrow \false$ — Z3 responde \texttt{sat} con modelo $x=5 \Rightarrow$ \textbf{[BUG]}.

\subsection{Suma con invariante de bucle}

\begin{lstlisting}[style=java, caption={Suma 1..5 con invariante completo}]
@post result >= 0
public static int sumLoop() {
    int i; int s;
    i = 1; s = 0;
    while (i <= 5)
    @invariant s >= 0 && i >= 1
    @variant   6 - i
    {
        s = s + i;
        i = i + 1;
    }
    return s;
}
\end{lstlisting}

\begin{table}[H]
  \centering
  \caption{VCs generadas para \texttt{sumLoop} y resultados de Z3}
  \label{tab:sumloop-vcs}
  \begin{tabular}{lp{7.5cm}l}
    \toprule
    \textbf{VC} & \textbf{Fórmula} & \textbf{Z3} \\
    \midrule
    Correctitud &
      $\true \Rightarrow (0\geq 0 \wedge 1\geq 1)$ &
      UNSAT \textcolor{green!60!black}{[OK]} \\
    Iniciación &
      $\true \Rightarrow (0\geq 0 \wedge 1\geq 1)$ &
      UNSAT \textcolor{green!60!black}{[OK]} \\
    Preservación &
      $(s\geq 0 \wedge i\geq 1 \wedge i\leq 5) \Rightarrow ((s+i)\geq 0 \wedge (i+1)\geq 1)$ &
      UNSAT \textcolor{green!60!black}{[OK]} \\
    Uso &
      $(s\geq 0 \wedge i\geq 1 \wedge i>5) \Rightarrow s\geq 0$ &
      UNSAT \textcolor{green!60!black}{[OK]} \\
    Terminación &
      $(s\geq 0 \wedge i\geq 1 \wedge i\leq 5) \Rightarrow (6-i)\geq 0$ &
      UNSAT \textcolor{green!60!black}{[OK]} \\
    \bottomrule
  \end{tabular}
\end{table}

\section{Nota sobre invariantes}

Un invariante debe ser \emph{suficientemente fuerte} para que las tres VCs sean
válidas. Un invariante débil como $s \geq 0$ no captura que $i \geq 1$, lo que
hace que Z3 encuentre contraejemplos con $i$ negativo para la VC de
preservación. La fortaleza del invariante es responsabilidad del programador;
el verificador comprueba que el invariante dado es \emph{correcto} (no que sea
\emph{completo}).
